In this section, we prove \Cref{thm:inaccurate_belief}, giving bounds on optimization of agents with inaccurate belief. The upper bound is in terms of the (weaker) absolute error and lower bounds are in terms of both absolute and relative error.

We use in this section the notion of total variational distance and coupling of random variables. For discussion of these topics, see \cite[Chapter 4.2]{levinmarkov}

\begin{repeatthm}{thm:inaccurate_belief}
Let $A$ be a perfect expected utility maximizer whose belief $\rho$ has absolute error $\epsilon$ with respect to the correct belief $\rho^*$. Then it is a $(\frac{2}{1-\gamma} - \frac{2}{1- \gamma(1-\epsilon)})$-optimizer with respect to $\rho^*$ and this bound is tight up to a factor of $2$. Moreover, if $\epsilon$ is the relative error, this bound is tight up to a factor of 4.
\end{repeatthm}

Before proving this, we will need the following lemma:
\begin{lemma} \label{lem:epsilon_bound_belief}
Let us consider the discounted expected utility with discount factor $\gamma$. Let $\rho$ have absolute error $\epsilon$ with respect to $\rho^*$. Then for any fixed policy $\pi$ and history $\mae_{<t}$.

\begin{gather}
|V_\rho^\pi(\mae_{<t}) - V_{\rho^*}^\pi(\mae_{<t})| \leq \frac{1}{1-\gamma} - \frac{1}{1- \gamma(1-\epsilon)}
\end{gather}

\end{lemma}
Before proving this, we will in turn need the \Cref{lem:belief_v_bound}; its proof is deferred to below.
\begin{lemma} \label{lem:belief_v_bound}
Let $\rho$ have absolute error $\epsilon$ with respect to $\rho^*$. Let $\mathcal{L}_\rho$ be the distribution on $(\mathcal{A}\times \mathcal{E})^*$ induced by $\rho$ and let
\[
\epsilon_{t} = \|\mathcal{L}_\rho (\mae_{k+t} | \mae_k) - \mathcal{L}_{\rho^*} (\mae_{k+t} | \mae_k) \|_{TV} = \sup_{\substack{S \subseteq (\mathcal{A \times E})^{t+k} \\ \text{s.t. first $k$ steps are}\\\text{equal to $\mae_{<k}$}}} |P_\rho(S | \mae_{<k}) - P_{\rho^*}(S | \mae_{<k})|
\]

Then it holds that
\[
\epsilon_{t} = 1-(1-\epsilon)^t
\]
\end{lemma}

\begin{proof}[Proof of \Cref{lem:epsilon_bound_belief}]
Let us bound
\begin{gather}
|V_{\rho}^{\pi}(\mae_{<k}) - V_{\rho^*}^\pi(\mae_{<k})| = |\mathbb{E}_{\rho}(\sum_{t=1}^\infty \gamma^{t-1} u(\mae_{<t+k}) | \mae_{<k}) - \mathbb{E}_{\rho^*}(\sum_{t=1}^\infty \gamma^{t-1} u(\mae_{<t+k}) | \mae_{<k})| \leq \\
\leq \sum_{t=1}^\infty \gamma^{t-1} |\mathbb{E}_\rho (u(\mae_{<t+k}) | \mae_{<k}) - \mathbb{E}_{\rho^*} (u(\mae_{<t+k}) | \mae_{<k})| \stackrel{(1)}{\leq} \sum_{t=1}^\infty \gamma^{t-1} \epsilon_t^\text{abs}  \leq \sum_{t=1}^\infty \gamma^{t-1}(1-(1-\epsilon)^t) = \\
= \frac{1}{\gamma}(\sum_{t=1}^\infty \gamma^{t} - \sum_{t=1}^\infty (\gamma(1-\epsilon))^t) = \frac{1}{1-\gamma} - \frac{1}{1- \gamma(1-\epsilon)}
\end{gather}

where, again, the linearity of expectation in infinite sums holds by the monotone convergence theorem and it, therefore, remains to argue inequality (1). By \Cref{lem:belief_v_bound}, it holds that the total variational distance is at most $\epsilon_t$ and, therefore, there exists a coupling of $X \sim \rho$ and $Y \sim \rho^*$ such that $P(X \neq Y) \leq \epsilon_t$. We can, therefore write for any function $f: (\mathcal{A \times E})^* \rightarrow [0,1]$
\begin{align}
E_\rho f &= E f(X) = E(f(X) | X = Y)P(X = Y) + E(f(X) | X \neq Y)P(X \neq Y) \\
E_{\rho^*} f &= E f(Y) = E(f(Y) | X = Y)P(X = Y) + E(f(Y) | X \neq Y)P(X \neq Y)
\end{align}
which we can now use to bound
\begin{align}
|E_\rho f - E_{\rho^*} f| =& |E(f(X) | X = Y)P(X = Y) + E(f(X) | X \neq Y)P(X \neq Y) \\&- E(f(Y) | X = Y)P(X = Y) + E(f(Y) | X \neq Y)P(X \neq Y)|\\
=& |E(f(X) | X \neq Y)P(X \neq Y) - E(f(Y) | X \neq Y)P(X \neq Y)|\\
\leq& P(X \neq Y) \leq \epsilon_t
\end{align}
where the first of the two inequalities holds because $f(X), f(Y) \in [0,1]$.

%holds because the probability of each elementary event at step $t$ differs between $\rho$ and $\rho^*$ by at most $\epsilon_t$, so the difference of the expectation is at most that. \jakub{give a better reason, this is wrong because it is not useing that utility is in $[0,1]$}

\end{proof}

\begin{proof}[Proof of \Cref{lem:belief_v_bound}]
Given a history $\mae_{<k+i}$, define the set $S_{k+i+1}(\mae{k+i})$ as the set of histories of length $k+i+1$ that are equal to $\mae_{<k+i}$ on the first $k+i$ timesteps. Now consider the probability distribution $\rho_{k+i+1, \mae_{<k+i}}$ as the distribution induced by $\rho$ on $S_{k+i+1}(\mae_{k+i})$ by conditioning on the first $k+i$ timesteps being equal to $\mae_{k+i}$ and similarly defined $\rho^*_{k+i+1, \mae_{k+i}}$.

By the \Cref{def:inaccurate_belief_abs}, we have for every $i \leq t$ that $\|\rho_{k+i+1, \mae_{<k+i}} - \rho^*_{k+i+1, \mae_{k+i}}\|_{TV} \leq \epsilon$. There, therefore, exist couplings $X_i \sim \rho_{k+i+1, \mae_{<k+i}}, Y_i \sim \rho^*_{k+i+1, \mae_{<k+i}}$ such that $P(X_i \neq Y_i)$. Let us define $X = (\mae_{<k}, X_1, \cdots X_t)$ and $Y = (\mae_{<k}, Y_1, \cdots, Y_t)$. These two sequences have the same distribution as $\mae_{<k+t}$ with respect to $\rho$ and $\rho^*$, respectively. Moreover, it holds that $P(X \neq Y) \leq 1- (1- \epsilon)^t$. We, therefore, have a coupling of $\mae_{<k+t}$ with respect to the two probbility measures $\rho, \rho^*$ and it follows that their total variational distance is at most $P(X \neq Y) \leq 1- (1- \epsilon)^t$.

%Let $\rho^*(e_k) = p_k$ and $\rho(e_k) = p_k + \epsilon_k$. Then 
%\[
%\rho(\mae_{<k+t} | \mae_{<k}) - \rho^*(\mae_{<k+t} | \mae_{<k}) = \prod_{i=k}^{k+t-1} (p_i + \epsilon_i) - \prod_{i=k}^{k+t-1} p_i \eqdef f(\vec{p}, \vec{\epsilon})
%\]
%We will now show that the points which maximise $|f(\vec{p}, \vec{\epsilon})|$ subject to constraints $|\epsilon_i| \leq \epsilon$, $0 \leq p_1 \leq 1$ and $0 \leq p_1+\epsilon_i \leq 1$ are $\vec{x_\text{1}} \eqdef (\overbrace{1-\epsilon, \cdots, 1-\epsilon}^{t\text{ times}}, \overbrace{\epsilon, \cdots, \epsilon}^{t\text{ times}})$ and $\vec{x_\text{2}} \eqdef (\overbrace{1, \cdots, 1}^{t\text{ times}}, \overbrace{-\epsilon, \cdots, -\epsilon}^{t\text{ times}})$.
%
%In order to maximise $|f|$, we need to either maximise or minimise  $f$. Let us first find the maximum of $f$. If $\epsilon_j \neq \epsilon$ for some $j$, we can increase the value of $f$ by increasing $\epsilon_j$. If this is not possible due to the constraint $p_j+\epsilon_j \leq 1$, we can simultaneously increase $\epsilon_j$ and decrease $p_j$ by the same amount, keeping $p_j+\epsilon_j$ constant. This also increases the value of $f$. Therefore, at the maximum of $f$, all $\epsilon_i = \epsilon$. We can further show that all $p_i = 1-\epsilon$ by evaluating $\partial f / \partial p_i$:
%
%\[
%\frac{\partial f}{\partial p_i} = \frac{\partial}{\partial p_i} \prod_j (p_j+\epsilon) - \frac{\partial}{\partial p_i} \prod_j p_j = \prod_{j\neq i} (p_j+\epsilon) -  \prod_{j\neq i} p_j
%\]
%This is always a positive number, so to maximise $f$, we need to choose the highest possible $p_j$ subject to the constraints, which is $p_j = 1 - \epsilon$. The point that maximises $f$ is therefore $x_\text{1} \eqdef (\overbrace{1-\epsilon, \cdots, 1-\epsilon}^{t\text{ times}}, \overbrace{\epsilon, \cdots, \epsilon}^{t\text{ times}})$.
%
%It can be shown by a similar argument that $\vec{x_\text{2}} \eqdef (\overbrace{1, \cdots, 1}^{t\text{ times}}, \overbrace{-\epsilon, \cdots, -\epsilon}^{t\text{ times}})$ minimises $f$. At both points, $|f(\vec{x_\text{1}})| = |f(\vec{x_\text{2}})| = 1-(1-\epsilon)^t$ and thus $\epsilon_t = 1 - (1-\epsilon)^t$.
%
%Note that the set of feasible solutions according to these constraints is exactly the set of possible probabilities according to $\rho^*$ and errors made by $\rho$ such that $\rho$ has absolute error of $\epsilon$ with respect to $\rho^*$. \jakub{say more}
%
\end{proof}


\begin{proof}[Proof of \Cref{thm:inaccurate_belief}]
The first part follows from \Cref{lem:epsilon_bound_belief} together with \Cref{lem:optimization_bound}.

We show tightness by constructing such utility, and belief. Both the action and perception sets are $\{0,1\}$. The utility is $1$ is all past percepts were $1$ and $0$ otherwise. The belief $\rho^*$ is such that when the last action performed was $1$, the perception is $1$ with probability one. If it was $0$, the perception is $1$ with probability $p_1 = 1-2\epsilon$ and $0$ otherwise. In belief $\rho$, independently of the action, the probability of $1$ is $p_2 = 1-\epsilon$. We can now assume that action $0$ is always taken. Then the discounted expected utility at step $k$ is $\gamma^{k-1} (1-2\epsilon)^k$, making the total expected discounted utility $\frac{1}{1- \gamma(1-2\epsilon)}$ instead of the optimal $\frac{2}{1-\gamma}$. The utility lost is then
\begin{gather}
\frac{1}{1-\gamma} - \frac{1}{1- \gamma(1-2\epsilon)} \leq \frac{1}{1-\gamma} - \frac{1}{1- \gamma(1-\epsilon)}
\end{gather}

In the relative case, the construction is the same except that we set $p_1 = \frac{1}{(1+\epsilon)^2}$ and $p_2 = \frac{1}{1+\epsilon}$. The expected discounted utility at step $t$ is then $\gamma^{t-1}\frac{1}{(1+\epsilon)^{2t}}$ resulting in this much utility lost:
\[
\frac{1}{1-\gamma} - \frac{1}{1- \frac{\gamma}{(1+\epsilon)^2}} \leq \frac{1}{1-\gamma} - \frac{1}{1- \frac{\gamma}{(1+\epsilon)}}
\]

Simplifying the ratio of this and $\frac{1}{1-\gamma} - \frac{1}{1- \gamma(1-\epsilon)}$ gives us (we skip the calculations)
\[
\frac{-\gamma +\epsilon +1}{\gamma  (\epsilon -1)+1}
\]
Simplifying the gradient of this, we get (we skip the calculations)
\[
\nabla_{\epsilon, \gamma} \left(\frac{-\gamma +\epsilon +1}{\gamma  (\epsilon -1)+1}\right) = \left(\frac{(\gamma -1)^2}{(\gamma  (\epsilon -1)+1)^2},-\frac{\epsilon ^2}{(\gamma  (\epsilon -1)+1)^2}\right)
\]
The function is, therefore, non-decreasing in $\epsilon$ and non-increasing in $\gamma$ and it has its maximum, subject to $0 \leq \epsilon, \gamma \leq 1$, at $\epsilon=1, \gamma=0$. Evaluating at this point, we get 2 as the value, meaning that the bound is at most a factor of 4 greater than the lower bound in the case of relative error. 
\end{proof}

In the next theorem, we show a lower bound in the average-case scenario which is asymtotically the same as the worst-case one. First we will need the following lemma:

\begin{lemma} \label{lem:probabilistic_bound}
Let $f: \mathbb{R}^n \rightarrow \mathbb{R}$ be a coordinate-wise non-decreasing function and let $X_1, \cdots, X_n$ be random variables. Let us for every $i \in [n]$ have $k_i$ random variables $X_i^1, \cdots, X_i^{k_i}$ such that $L_{X_i^j} \simeq X_i$. Let $f_l = f(X_1^{j_{1,l}}, \cdots, X_n^{j_{n,l}})$ for $l \in [L]$ for some $L$. Consider the set of random variables $\mathcal{F} = \{f_l: l \in [L]\}$ for some set of $j_{i,l}$ s.t. $j_{i,l} \in [k_i]$. Let us consider $f_\text{max} \eqdef \max \mathcal{F}$. Finally, let $l_\text{max}$ be such that $f_\text{max} = f(X_1^{j_{1,l_\text{max}}}, \cdots, X_n^{j_{n,l_\text{max}}})$.

Then $P(X_i^{j_{i,l_\text{max}}} \leq x) \leq P(X_i \leq x)$ for all $i \in [n]$ and $x$.
\end{lemma}
The interpretation is that if we have a non-decreasing function (possibly in multiple variables) and evaluate it several times on possibly different instances of the same random variables, and we take the arguments that resulted in the highest value of the function, then the arguments are generally speaking larger than the value of the respective random variables. More precisely, taking the arguments which resulted in $f_\text{max}$, there exists a coupling of the $i$-th argument with $X_i$ such that the argument is always greater than $X_i$.
\begin{proof}
We have $f_\text{max} \simeq f_1 | f_1 \geq f_i, i \in \{2,3,\cdots,n\}$. We prove that for any value $C$, it holds that $P(X_i^{j_{i,1}} \leq x | f_1 \geq C) \leq P(X_i \leq x)$ and then set $C = \max(f_i, i \in \{2,3,\cdots,n\})$.

Let $p_i(x) = P(X_i=x)$. We consider the set $T$ of $x_1, \cdots, x_n$ such that $f(x_1, \cdots, x_n) \geq C$. Let $S_i(x_i) = \{\mathbf{x} = (x_1, \cdots, x_{i-1}, x_{i}, x_{i+1}, \cdots, x_n), s.t. f(\mathbf{x}) \geq C\}$ and $W_i(x)$ be the hyperplane with the $i$-th component equal to $x$. In other words, $S_i(x) = T \cap W_i(x)$. Note that $S_i(x)$ is non-decreasing. We have that the probabilities
\begin{gather}
p_i(x_i) = \sum_{x\in W_i(x)} \prod_{j=1}^n p_j(x) \\
p_i(x_i | f(x_1, \cdots, x_{i-1}, x_{i}, x_{i+1}, \cdots, x_n) \geq C) = c\sum_{x\in S_i(x)} \prod_{j=1}^n p_j(x)
\end{gather}
for $c = P(f(x) \geq C)^{-1}$ by the Bayes' theorem. Now, because $S_i(x)$ is non-decreasing, we have that for every $i \in [n]$
\[
r_i(x_i) \eqdef \frac{p(x_i | f \geq C)}{p(x_i)} = \frac{c\sum_{x\in S_i(x)} \prod_{j=1}^n p_j(x)}{\sum_{x\in W_i(x)} \prod_{j=1}^n p_j(x)}
\]
is non-decreasing. Since $r_i(x)$ is non-decreasing there exists $x$ (we allow $x$ to be infinite, even though it can be shown that this is never the case) such that for any $x' \leq x$, it holds that $p_i(x') \geq p_i(x' | f \geq C)$ and $p_i(x') \leq p_i(x' | f \geq C)$ for $x' \geq x$. Assume that $x' \geq x$ as the argument for the other direction is analogous and slightly simpler. We have for $C$ equal to 
\begin{gather}
P(X_i^{j_{i,l_\text{max}}} \leq x) = 1-P(X_i^{j_{i,l_\text{max}}} > x) = 1-\sum_{y:y>x} p_i(y | f \geq C) \leq \\ \leq 1-\sum_{y:y>x} p_i(y) = 1-P(X_i > x) = P(X_i \leq x)
\end{gather}
\end{proof}

\begin{repeatthm}{thm:avg_case_belief}
	Let $A$ be a perfect expected utility maximizer whose belief $\rho$ is such that for any $e_t,\mae_{<t}a_t$, the value $\rho(e_t|\mae_{<t}a_t)$ is chosen randomly from the set $\{\max(0,\rho^*(e_t|\mae_{<t}a_t)-\epsilon), \min(1,\rho^*(e_t|\mae_{<t}a_t)+\epsilon)\}$ in the case of absolute error and $\{\frac{\rho^*(e_t|\mae_{<t}a_t)}{1+\epsilon}, \min(1,(1+\epsilon)\rho^*(e_t|\mae_{<t}a_t))\}$ in the case of relative error.
	
	Then, in expectation, the amount of expected discounted utility lost is respectively
	\begin{gather}
	\frac{1}{1-\gamma} - \frac{1}{1- \gamma(1-\epsilon/8)}\\
	\frac{1}{1-\gamma} - \frac{1}{1- \gamma(1-\epsilon/16)}
	\end{gather}
	
	and this is equal to the upper bound for non-random error up to a factor of $16$ and $32$, respectively.
\end{repeatthm}
\begin{proof}
We use a similar construction to the lower bound in the worst case. Both the action and perception sets are $\{0,1\}$. The utility is $1$ if all past percepts were $1$ and $0$ otherwise. The belief $\rho^*$ is such that when the last action performed was $1$, the perception is $1$ with probability $1$. If the last action was $0$, the perception is $1$ with probability $1-\epsilon$ and $0$ otherwise. We call these the $1$-path and $0$-path (note the asymetry in the definitions).

Assume that we are in some vertex of the tree and that the believed probability of perception 1 after action 0 is 1. The probability that the subtree of action 0 has optimum greater or equal to that of the subtree of action 1 happens by symmetry with probability at least $1/2$. Moreover, if the error of action $0$ is in the positive direction, we choose action $0$ with probability at least $1/2$ giving us probability at least $1/8$ of performing action 0.

Therefore, we need to lowerbound the probability that the error is in the positive direction --- in any fixed vertex, this happens with probability exactly $1/2$ but in our case, the probability is conditioned on the fact that the optimal path goes through the vertex. For this, we use \Cref{lem:probabilistic_bound}. For any fixed path, the expected discounted utility is non-decreasing in the absolute errors. We can, therefore, use the lemma to show that if we choose the optimal path, the probability of any of the errors being in the positive direction is at least that for any fixed action (that is $1/2$). Therefore, in each step, we have probability at least $1/8$ that the agent performs action $0$.

We have $\mathbb{E}((1-\epsilon)^{X_t}) \leq \frac{7}{8} + \frac{1}{8}(1-\epsilon) = 1-\frac{\epsilon}{8}$ and by independence $\mathbb{E}\epsilon_t = \mathbb{E}(1-\prod_{i=1}^t (1-\epsilon)^{X_i}) \geq 1-(1-\frac{1}{8}\epsilon)^t$. By a similar argument as above, this gives the lost expected discounted utility of 

\[
\epsilon'_\text{abs} \eqdef \frac{1}{1-\gamma} - \frac{1}{1- \gamma(1-\epsilon/8)}
\]

Similarly for the relative error, $\mathbb{E}(\frac{1}{1+\epsilon}^{X_t}) \leq \mathbb{E}((1-\frac{1}{2}\epsilon)^{X_t}) \leq 1-\frac{1}{16}\epsilon$ and $\mathbb{E}\epsilon_t = \mathbb{E}(1-\prod_{i=1}^t \frac{1}{1+\epsilon}^{X_i}) \geq 1-(1-\frac{1}{16}\epsilon)^t$. Then the expectation of the expected discounted utility lost is no more than
\[
\epsilon'_\text{rel} \eqdef \frac{1}{1-\gamma} - \frac{1}{1- \gamma(1-\epsilon/16)}
\]

Now we bound the ratio of the upper bound $f_\text{bel}(\epsilon, \gamma)$ and the lower bound $\epsilon'_\text{abs}$ (we skip the simplification)

\[
\frac{f_\text{bel}(\epsilon, \gamma)}{\epsilon'_\text{abs}} = \frac{2 (\gamma  (\epsilon -8)+8)}{\gamma  (\epsilon -1)+1} \leq  \frac{2 (\gamma  (8\epsilon -8)+8)}{\gamma  (\epsilon -1)+1} = 16
\]
Similarly for the relative error, we get an upper bound of $32$ on the ratio.

\end{proof}

